\section{Discussion}\label{sec:discussion}

The size of a fitted entropy production says little by itself. The fitted OU model is misspecified---a windowed state is not Markov---but it is the model that windowed early-warning pipelines fit, and Proposition~\ref{prop:factor} separates its current, measured by $\omega_{\rm fit}$, from the innovation geometry, measured by $\mathcal G$. The latter can vary by orders of magnitude between states with similar currents: for the Gaussian observer with noise-free windows $\mathcal G$ exceeds $10^4$, whereas in the market forecasting state it is \GPca, and re-fitting a reduction can change $\dot e_p$ several-fold without any change in the dynamics (Sec.~\ref{sec:pca}). Reporting $\omega$ alongside $\dot e_p$ makes such numbers comparable.

Irreversibility measured on windows of one observable is a property of the windows (Corollary~\ref{cor:single}). Second-order indicators built from several trailing windows of a single signal cannot distinguish a reversible from an irreversible system; detecting the latter requires higher-order statistics or genuinely distinct observed channels. For multichannel states the time-reversed data provide a model-free baseline, because the observer term is what survives reversal, and centered windows provide a second one where the features allow it. When irreversibility is used as an early-warning signal \cite{kooloth2026,xu2023} on states assembled from trailing features, the observer term should be removed first; Eq.~\eqref{eq:estimator} costs one extra pass of the feature pipeline.

In financial data the two terms have distinct origins. The observer term lives in comparisons of volatility over different windows and is present at full strength in surrogates whose volatility dynamics is reversible. The physical term is the odd part of the cross-covariance between signed returns and later volatility, the leverage effect; it survives centered windows and disappears with random signs. Higher-order asymmetry of volatility itself, the Zumbach effect \cite{zumbach2009,lynch2003}, exists in the same data (Sec.~\ref{sec:robust}) but is invisible to second-order window statistics. This is consistent with the finding that time-reversal violations of returns measured through feedback kernels reflect the ARCH mechanism rather than leverage \cite{chicheportiche2014}: that statistic is of higher order, ours is of second order, and the two probe different channels.

For information loss the lesson is that the relevant clock is the slowest relaxation, not the local one. In the double well most of the oracle's variance lies between the wells and decorrelates at the Kramers rate, so trailing windows much longer than the local relaxation time still preserve most of the transition information; a single-rate theory based on local curvature overestimates the loss by a factor of \SingleFacLo to \SingleFacHi. In the market the windows themselves are not the main loss: averaging the latent variance, the 5- and 20-day windows would lose \EllMktPair of the information about the next volatility excursion, but they average squared returns and lose \EllMktRV (Sec.~\ref{sec:mktinfo}). Forecasts that integrate a fitted drift, such as committor-type hitting probabilities, inherit both effects of the observer; a companion paper evaluates such forecasts on the same market state \cite{xing2026b}.

The study has limits. Proposition~\ref{prop:decomp} is exact for linear filters and second-order statistics; for nonlinear features such as running maxima the split is approximate and the centering control fails, and higher-order irreversibility is outside its scope. The empirical part rests on one equity index and two cryptocurrencies at daily sampling. The stochastic-volatility check uses one- and two-factor models, which bracket but do not pin down the observed values, and the market reading of the information loss rests on the one-factor calibration. The two-rate approximation for the double well captures the mechanism but underestimates the loss at the widest windows by up to \TwoErrAbs.

\section{Conclusion}\label{sec:conclusion}

Trailing windows add a current of their own to whatever current a system carries; at second order the two separate exactly under time reversal, and the same windows destroy transition information at rates set by the system's relaxation, including Kramers switching in metastable systems. For an Ornstein--Uhlenbeck input both effects are in closed form, and the large fitted dissipation of reversible inputs is a geometric consequence of collinear innovations. In daily S\&P 500 data the volatility windows carry the observer's current and the leverage effect carries the physical one, and the forecasting state lacks about half of the information that the latent volatility carries about its next excursion. The same accounting applies to any early-warning state built from windows, for instance in climate and ecological indicators.
